% !TeX root = ../main.tex
% !TeX encoding = utf8
	\section{XÂY DỰNG KHOẢNG TIN CẬY NORMAL CI VÀ BASIC CI}
	
	\subsection{Khoảng tin cậy Bootstrap chuẩn tắc (Normal CI)}
	
	\subsubsection{Cơ sở và giả định}
	Normal Bootstrap CI dùng sai số chuẩn Bootstrap trong công thức khoảng dựa trên xấp xỉ Chuẩn \cite{buithuytrang2026}. Phương pháp giả định phân phối lấy mẫu của $\hat{\theta}_n$ xấp xỉ Chuẩn quanh tham số $\theta$ \cite{efron1993}:
	\begin{equation}
		\hat{\theta}_n \approx N\left(\theta, se(\hat{\theta}_n)^2\right) \quad \text{hoặc} \quad \frac{\hat{\theta}_n - \theta}{se(\hat{\theta}_n)} \approx N(0, 1)
	\end{equation}
	Đặt
	\[
	Z=\frac{\hat{\theta}_n-\theta}{se(\hat{\theta}_n)}.
	\]
	Nếu $Z$ xấp xỉ phân phối chuẩn tắc, thì
	\[
	\mathbb{P}\left(
	-z_{1-\alpha/2}\leq Z\leq z_{1-\alpha/2}
	\right)\approx 1-\alpha.
	\]
	Giải bất đẳng thức theo $\theta$ cho ta khoảng
	$\hat{\theta}_n\pm z_{1-\alpha/2}se(\hat{\theta}_n)$. Normal Bootstrap CI
	thay sai số chuẩn chưa biết bằng sai số chuẩn Bootstrap $\widehat{se}^*$.
	Thay vì tính toán $se(\hat{\theta}_n)$ bằng các công thức giải tích phức tạp, Normal CI sử dụng ước lượng sai số chuẩn Bootstrap $\widehat{se}^*$ thu được từ mô phỏng Monte Carlo \cite{efron1993}.
	
	\subsubsection{Công thức xây dựng khoảng tin cậy}
	Khoảng tin cậy Bootstrap Chuẩn tắc ở độ tin cậy $1 - \alpha$ được xác định bởi công thức \cite{buithuytrang2026}:
	\begin{equation}
		CI_{\text{Normal}} = \left[ \hat{\theta}_n - z_{1-\alpha/2} \cdot \widehat{se}^*, \quad \hat{\theta}_n + z_{1-\alpha/2} \cdot \widehat{se}^* \right]
	\end{equation}
	trong đó:
	\begin{itemize}
		\item $\hat{\theta}_n = t(x)$ là ước lượng điểm tính trên mẫu dữ liệu gốc \cite{buithuytrang2026}.
		\item $z_{1-\alpha/2} = \Phi^{-1}(1 - \alpha/2)$ là phân vị mức $1-\alpha/2$ của phân phối Chuẩn tắc $N(0,1)$ \cite{buithuytrang2026}. Với $\alpha=0.05$, giá trị này xấp xỉ $1.96$.
		\item $\widehat{se}^* = \sqrt{\frac{1}{B-1}\sum_{b=1}^{B} \left(\hat{\theta}^{*(b)} - \bar{\theta}^*\right)^2}$ là độ lệch chuẩn mẫu của $B$ bản sao Bootstrap \cite{buithuytrang2026}.
	\end{itemize}
	
	\subsection{Khoảng tin cậy Bootstrap cơ bản (Basic CI)}
	
	\subsubsection{Cơ sở và phép đảo phân vị (quantile flipping)}
	Basic CI không yêu cầu phân phối lấy mẫu của $\hat{\theta}_n$ là Chuẩn đối xứng \cite{efron1993}. Phương pháp dùng phân phối Bootstrap của sai số $\delta^*=\hat{\theta}_n^*-\hat{\theta}_n$ để xấp xỉ phân phối lấy mẫu của $\delta=\hat{\theta}_n-\theta$ \cite{davison1997}:
	\begin{equation}
		\mathbb{P}\left(\hat{\theta}_n - \theta \le t\right) \approx \mathbb{P}^*\left(\hat{\theta}_n^* - \hat{\theta}_n \le t\right)
	\end{equation}
	Ở đây, $\mathbb{P}^*$ là xác suất trong thế giới Bootstrap khi giữ cố định
	mẫu quan sát ban đầu. Dấu xấp xỉ là giả định then chốt của Basic CI:
	phân phối sai số Bootstrap được dùng để xấp xỉ phân phối sai số lấy mẫu.
	Các phép biến đổi trong định lý sau đó là đại số dựa trên giả định này;
	chúng không tự chứng minh rằng phép xấp xỉ luôn chính xác.
	\begin{theorem}[\textbf{Chứng minh Đại số Cơ chế Quantile Flipping}]
		Gọi $q_{\alpha/2}^*$ và $q_{1-\alpha/2}^*$ lần lượt là các phân vị mức $\alpha/2$ và $1-\alpha/2$ của tập hợp các bản sao Bootstrap $\{\hat{\theta}^{*(1)}, \dots, \hat{\theta}^{*(B)}\}$ \cite{davison1997}. Theo định nghĩa phân vị thực nghiệm:
		\begin{equation}
			\mathbb{P}^*\left( q_{\alpha/2}^* \le \hat{\theta}_n^* \le q_{1-\alpha/2}^* \right) \approx 1 - \alpha
		\end{equation}
		Biến đổi biểu thức trong dấu xác suất bằng cách trừ $\hat{\theta}_n$ ở tất cả các vế \cite{davison1997}:
		\begin{equation}
			\mathbb{P}^*\left( q_{\alpha/2}^* - \hat{\theta}_n \le \hat{\theta}_n^* - \hat{\theta}_n \le q_{1-\alpha/2}^* - \hat{\theta}_n \right) \approx 1 - \alpha
		\end{equation}
		Áp dụng giả định xấp xỉ độ lệch, ta thay $\hat{\theta}_n^* - \hat{\theta}_n$ bằng $\hat{\theta}_n - \theta$ \cite{davison1997}:
		\begin{equation}
			\mathbb{P}\left( q_{\alpha/2}^* - \hat{\theta}_n \le \hat{\theta}_n - \theta \le q_{1-\alpha/2}^* - \hat{\theta}_n \right) \approx 1 - \alpha
		\end{equation}
		Giải bất phương trình trên đối với tham số tổng thể thực $\theta$ \cite{davison1997}:
		\begin{align*}
			& q_{\alpha/2}^* - \hat{\theta}_n \le \hat{\theta}_n - \theta \le q_{1-\alpha/2}^* - \hat{\theta}_n \\
			\iff & q_{\alpha/2}^* - 2\hat{\theta}_n \le -\theta \le q_{1-\alpha/2}^* - 2\hat{\theta}_n \\
			\iff & 2\hat{\theta}_n - q_{1-\alpha/2}^* \le \theta \le 2\hat{\theta}_n - q_{\alpha/2}^* \quad (\text{đpcm})
		\end{align*}
	\end{theorem}
	
	\subsubsection{Công thức xây dựng khoảng tin cậy}
	Khoảng tin cậy Bootstrap Cơ bản ở độ tin cậy $1 - \alpha$ có dạng \cite{buithuytrang2026}:
	\begin{equation}
		CI_{\text{Basic}} = \left[ 2\hat{\theta}_n - q_{1-\alpha/2}^*, \quad 2\hat{\theta}_n - q_{\alpha/2}^* \right]
	\end{equation}
	\subsection{Thuật toán xây dựng hai khoảng tin cậy}
	\begin{algorithm}[H]
		\Input{Mẫu dữ liệu $x = (x_1, \dots, x_n)$; Số lần lấy mẫu lại $B$; Mức ý nghĩa $\alpha$ \cite{buithuytrang2026}.}
		\Output{Khoảng tin cậy $CI_{\text{Normal}}$ và $CI_{\text{Basic}}$ ở độ tin cậy $(1-\alpha) \cdot 100\%$ \cite{buithuytrang2026}.}
		\BlankLine
		\textbf{Bước 1: Tính ước lượng trên mẫu gốc}\;
		Tính $\hat{\theta}_n = t(x_1, \dots, x_n)$ \cite{buithuytrang2026}\;
		
		\textbf{Bước 2: Rút mẫu Bootstrap Monte Carlo}\;
		\For{$b = 1, 2, \dots, B$}{
			Rút ngẫu nhiên có hoàn lại $n$ phần tử từ $x$ để tạo $x^{*(b)}$ \cite{buithuytrang2026}\;
			Tính bản sao Bootstrap: $\hat{\theta}^{*(b)} = t(x^{*(b)})$ \cite{buithuytrang2026}\;
		}
		
		\textbf{Bước 3: Tính các đặc trưng phân phối Bootstrap}\;
		Sắp xếp $B$ bản sao theo thứ tự tăng dần: $\hat{\theta}^{*(1)} \le \hat{\theta}^{*(2)} \le \dots \le \hat{\theta}^{*(B)}$ \cite{buithuytrang2026}\;
		Tính sai số chuẩn Bootstrap $\widehat{se}^* = \sqrt{\frac{1}{B-1}\sum_{b=1}^{B} (\hat{\theta}^{*(b)} - \bar{\theta}^*)^2}$ \cite{buithuytrang2026}\;
		Với phân vị thực nghiệm type 1, lấy chỉ số $k_L=\max\{1,\lceil B\alpha/2\rceil\}$ và $k_U=\lceil B(1-\alpha/2)\rceil$; đặt $q_{\alpha/2}^*=\hat{\theta}^{*(k_L)}$ và $q_{1-\alpha/2}^*=\hat{\theta}^{*(k_U)}$ \cite{buithuytrang2026}\;
		
		\textbf{Bước 4: Xây dựng khoảng tin cậy}\;
		Tính $CI_{\text{Normal}} = [\hat{\theta}_n - z_{1-\alpha/2}\widehat{se}^*, \quad \hat{\theta}_n + z_{1-\alpha/2}\widehat{se}^*]$ \cite{buithuytrang2026}\;
		Tính $CI_{\text{Basic}} = [2\hat{\theta}_n - q_{1-\alpha/2}^*, \quad 2\hat{\theta}_n - q_{\alpha/2}^*]$ \cite{buithuytrang2026}\;
		\Return $CI_{\text{Normal}}$ và $CI_{\text{Basic}}$ \cite{buithuytrang2026}\;
		\caption{Thuật toán 4 bước xây dựng Khoảng tin cậy Normal CI và Basic CI}
	\end{algorithm}
	
	\subsection{Liên hệ với ví dụ tính toán}

Các bước tính số, mã R và kết quả minh họa được trình bày tập trung ở
Phần~\ref{sec:examples}. Cách tổ chức này giúp phần lý thuyết tập trung vào
thuật toán, còn các ví dụ và đầu ra R được đối chiếu tại cùng một vị trí.

